\section{The Angels in Liturgy, Devotion, and the Lives of the Saints}

Before a Catholic ever opens a treatise on separate substances, he has already spoken to the angels at Mass. He confesses his sins ``to blessed Michael the Archangel,'' prays that his voice be admitted among the Dominations, Powers, Virtues, and Seraphim, sings the Seraphic \emph{Sanctus}, and asks Almighty God to bid his holy angel carry the consecrated offering to the altar on high. Catholic angelology begins and ends in worship, where theology becomes prayer and the Church on earth joins the heavenly city already keeping liturgy around the Lamb.

\subsection{The Angels in the Eucharistic Liturgy}

Across the Roman Rite, the Eucharistic liturgy presupposes an unbroken communion between the earthly assembly and the celestial choirs. In the \emph{Confiteor} of the 1962 Roman Missal (and pre-1955 editions such as the 1889 Descl\'ee \emph{Missale Romanum}), the penitent confesses \emph{beato Michaeli Archangelo} (``to blessed Michael the Archangel'') immediately after the Blessed Virgin Mary and asks his intercession alongside the Apostles and all the Saints. At the threshold of the Eucharistic Prayer, the Preface summons the heavenly court by name. In the classic Roman Common and Lenten Prefaces, the priest prays: \emph{Per quem majestatem tuam laudant Angeli, adorant Dominationes, tremunt Potestates. Coeli, coelorumque Virtutes, ac beata Seraphim, socia exsultatione concelebrant. Cum quibus et nostras voces ut admitti jubeas, deprecamur, supplici confessione dicentes}---``through whom the Angels praise thy majesty, the Dominations adore, the Powers tremble; the Heavens and the Virtues of heaven and the blessed Seraphim celebrate with united exultation; with whom we beseech thee to bid our voices also be admitted, saying in humble confession.'' The assembly then joins Isaias's Seraphim in the \emph{Sanctus}: \emph{Sanctus, Sanctus, Sanctus Dominus Deus Sabaoth. Pleni sunt coeli, et terra gloria tua, Hosanna in excelsis} (``Holy, Holy, Holy, Lord God of hosts; heaven and earth are full of thy glory; Hosanna in the highest'').

After the Consecration, the Roman Canon bows deeply in the prayer \emph{Supplices te rogamus}:

\begin{quotation}
\emph{Supplices te rogamus, omnipotens Deus: jube haec perferri per manus sancti Angeli tui in sublime altare tuum, in conspectu divinae majestatis tuae: ut quotquot ex hac altaris participatione, sacrosanctum Filii tui Corpus et Sanguinem sumpserimus, omni benedictione coelesti et gratia repleamur. Per eundem Christum Dominum nostrum. Amen.}
\end{quotation}

As rendered in the 1945 Lasance \emph{New Roman Missal}: ``We humbly beseech Thee, almighty God, to command that these our offerings be borne by the hands of Thy holy angel, to Thine altar on high, in the presence of Thy divine majesty, that as many of us as shall receive the most sacred body and blood of Thy Son by partaking thereof from this altar may be filled with every heavenly blessing and grace. Through the same Christ our Lord. Amen.'' Echoing the angel with the golden censer in Apocalypse 8:3--4, the Canon sets the Church's earthly sacrifice directly inside the heavenly liturgy.

\subsection{The Liturgical Calendars}

Both the Western Roman calendars and the Byzantine Menaion appoint proper days to honor the holy angels. Comparing the pre-1955 Roman Missals, the 1962 Roman Missal, the postconciliar General Roman Calendar with the 1981 \emph{Ordo Lectionum Missae}, and the Byzantine Menaion shows how the Church's liturgical remembrance of the angels has been arranged across rites and editions:

\begin{itemize}
\item \textbf{24 March (St.~Gabriel the Archangel):} Extended to the universal Roman calendar in 1921 and kept in the 1962 Roman Missal as a Class~III feast (\emph{S.~Gabrielis Archangeli}), appointing Dan 9:21--26 and Luke 1:26--38. In the postconciliar General Roman Calendar, St.~Gabriel is joined with Sts.~Michael and Raphael on 29 September.
\item \textbf{8 May (Apparition of St.~Michael the Archangel):} Present in pre-1955 Roman Missals (1889 Descl\'ee, 1922 Tours Mame, and 1945 Lasance) as a Greater Double commemorating the sanctuary of Monte Gargano in Apulia, with Apoc 1:1--5, Matt 18:1--10, and the Paschal Alleluia \emph{Sancte Michael Archangele, defende nos in proelio: ut non pereamus in tremendo judicio} (``St.~Michael the Archangel, defend us in battle, that we may not perish in the dreadful judgment''). Suppressed in the 1960/1962 rubrical reform, it does not appear in either the 1962 Missal or the postconciliar calendar.
\item \textbf{29 September (St.~Michael / Sts.~Michael, Gabriel, and Raphael):} In the 1962 Roman Missal, this date is the Class~I feast \emph{In Dedicatione S.~Michaelis Archangeli}, reading Apoc 1:1--5 and Matt 18:1--10, with the Collect \emph{Deus, qui, miro ordine, Angelorum ministeria hominumque dispensas} (``O God, who in a wonderful order dost dispense the ministries of angels and of men''). In the postconciliar General Roman Calendar and 1981 \emph{Ordo Lectionum Missae} (n.~647), it is kept as the Feast of \emph{Saints Michael, Gabriel, and Raphael, Archangels}, reading Dan 7:9--10, 13--14 or Apoc 12:7--12a, Ps 137 [Heb 138], and John 1:47--51, while retaining the ancient Latin Collect \emph{Deus, qui, miro ordine} verbatim.
\item \textbf{2 October (The Holy Guardian Angels):} In the 1962 Roman Missal, a Class~III feast (\emph{Ss.~Angelorum Custodum}) appointing Ex 23:20--23, Ps 90:11--12 [Heb 91:11--12], and Matt 18:1--10, with the Collect \emph{Deus, qui ineffabili providentia sanctos Angelos tuos ad nostram custodiam mittere dignaris} (``O God, who in ineffable providence dost deign to send thy holy angels to guard us''). In the postconciliar calendar and 1981 \emph{Ordo Lectionum Missae} (n.~650), it is an Obligatory Memorial appointing Ex 23:20--23, Ps 90 [Heb 91], and the proper Gospel Matt 18:1--5, 10, retaining the Latin Collect \emph{Deus, qui ineffabili providentia} verbatim.
\item \textbf{24 October (St.~Raphael the Archangel):} Added to the universal calendar in 1921 and kept in the 1962 Roman Missal as a Class~III feast (\emph{S.~Raphaelis Archangeli}), appointing Tob 12:7--15 and John 5:1--4 (the angel stirring the pool of Bethsaida). In the postconciliar calendar, St.~Raphael is celebrated on 29 September, while 24 October is the Optional Memorial of St.~Anthony Mary Claret.
\item \textbf{8 November (Byzantine Synaxis of the Archangels):} In the Byzantine Rite, the primary universal feast is the \emph{Synaxis of the Archangels and All the Bodiless Powers} on 8 November---the eighth day of the ninth month from March, recalling the nine choirs and the eighth day of the age to come---kept in the Byzantine Catholic Archeparchy of Pittsburgh Menaion as a Feast with an All-night Vigil.
\end{itemize}

\subsection{Ecclesial Regulation of Angelic Devotion and Names}

Catholic worship draws an unyielding boundary between \emph{latria}---supreme adoration due to God alone---and \emph{dulia}, the subordinate veneration and invocation offered to God's holy angels and saints. St.~Paul warned the Colossians against false ``religion of angels'' detached from Christ the Head (Col 2:18--19), and twice in the Apocalypse the angel stops St.~John from falling at his feet: ``See thou do it not. I am thy fellow servant \dots\ Adore God'' (Apoc 19:10; 22:8--9). Explaining the First Commandment, the \emph{Roman Catechism} of 1566 clarifies that the angels refused only supreme divine worship, while the Church rightly honors them as God's ministers and loving intercessors: ``Their intercession, therefore, we invoke, because they always see the face of God, and are constituted by him the willing advocates of our salvation'' (Part~III, First Commandment). In its preface to the Lord's Prayer, the same Catechism compares our Heavenly Father to parents sending guardians along a robber-infested road, placing angels over each of us ``guarded by whose vigilant care and assistance, we may escape the ambushes of our enemies.''

Because unhealthy curiosity about invisible spirits readily slips into superstition, the Church has guarded angelic devotion since antiquity---above all by restricting proper names to those revealed in Sacred Scripture. Around 343--381, Canon~35 of the Council of Laodicea condemned secretive Phrygian angelolatry (recorded by Karl Joseph von Hefele, \emph{History of the Councils}, vol.~II): ``Christians shall not forsake the Church of God and turn to the worship of angels, thus introducing a cultus of the angels. This is forbidden. Whoever, therefore, shows an inclination to this hidden idolatry, let him be anathema, because he has forsaken our Lord Jesus Christ, the Son of God, and gone over to idolatry.'' The Greek canon uses the phrase \gk{onomazein angelous} (``naming'' or invoking angels), and four centuries later the Roman Synod of 745 under Pope St.~Zacharias applied the same discipline against the heterodox preacher Aldebert, condemning his prayer invoking eight unscriptural angels and ruling that the Church knows only three angelic names from Scripture: Michael, Gabriel, and Raphael (reaffirmed in Charlemagne's 789 capitulary).

The modern Magisterium maintains the exact same guardrails. On 6 June 1992, the Congregation for the Doctrine of the Faith issued its decree on the association \emph{Opus Angelorum} (\emph{AAS}~84, pp.~805--806), reiterating earlier 1983 norms that forbid spreading theories or invoking proper names of angels derived from private revelations (specifically those attributed to Gabriele Bitterlich), prohibit consecrations to the angels, and ban unauthorized liturgical insertions. On 17 December 2001, the Congregation for Divine Worship and the Discipline of the Sacraments echoed this rule in the \emph{Directory on Popular Piety and the Liturgy} (\S\S212--217), commending authentic liturgical and personal devotion to the holy angels and guardian angels while explicitly discouraging the practice of assigning proper names to angels beyond the three revealed in Sacred Scripture: Michael, Gabriel, and Raphael.

\subsection{Received Vocal Prayers to the Holy Angels}

Two brief prayers, rooted in medieval and papal tradition and confirmed in public-domain Catholic prayer books, anchor daily Catholic invocation of the guardian angel and the Prince of the heavenly host.

\begin{studybox}{\emph{Angele Dei} (1910 \emph{Raccolta}, No.~296)}
\emph{Angele Dei, qui custos es mei, me tibi commissum pietate superna illumina, custodi, rege, et guberna. Amen.}\par\medskip
Angel of God, whom God hath appointed to be my guardian, enlighten and protect, direct and govern me. Amen.\par\medskip
Angel of God, my guardian dear,\\
To whom his love commits me here,\\
Ever this day be at my side,\\
To light and guard, to rule and guide. Amen.
\end{studybox}

\begin{studybox}{Prayer to St.~Michael the Archangel (Pope Leo XIII; 1945 Lasance \emph{New Roman Missal}, p.~797)}
\emph{Sancte Michael Archangele, defende nos in praelio; contra nequitiam et insidias diaboli esto praesidium. Imperet illi Deus, supplices deprecamur: tuque, Princeps militiae caelestis, Satanam aliosque spiritus malignos, qui ad perditionem animarum pervagantur in mundo, divina virtute in infernum detrude. Amen.}\par\medskip
St.~Michael, the archangel, defend us in battle, be our protection against the malice and snares of the devil. We humbly beseech God to command him, and do thou, O prince of the heavenly host, by the divine power thrust into hell Satan and the other evil spirits who roam through the world seeking the ruin of souls. Amen.
\end{studybox}

Alongside this short Leonine prayer recited after Low Mass, Pope Leo XIII also promulgated a longer prayer to St.~Michael by \emph{Motu Proprio} on 25 September 1888 (1910 \emph{Raccolta}, No.~292) and indulgenced the traditional liturgical antiphon \emph{Sancte Michael Archangele, defende nos in praelio, ut non pereamus in tremendo iudicio} on 19 August 1893 (\emph{Acta Sanctae Sedis}, vol.~26, p.~183)---the very Alleluia verse sung on St.~Michael's feast in the Roman Missal.

\subsection{The Saints' Devotion to the Angels}

The lives and writings of the canonized saints---belonging to \textbf{pious tradition}---show how deeply the doctrine of the angels shapes Christian conscience, contemplation, and trust.

Preaching on Psalm 90 (\emph{Sermons on ``Qui habitat,''} Sermon~XII), St.~Bernard of Clairvaux drew from verse~11 a classic three-part rule of life long read in the Divine Office for October~2:

\begin{quotation}
``He hath given His angels charge over thee: to keep thee in all thy ways.'' O my brother, with how much reverence, with how much gratitude, with how much confidence ought not these words to inspire thee! With how much reverence for a presence so august! With how much gratitude for benevolence so great! With how much confidence in a keeping so secure! Walk therefore with caution, remembering that the angels of God accompany thee in all thy ways, as the Lord has commanded them. In every place, whether public or private, show respect for thy angel. Surely thou wouldst not dare to do in his presence what thou wouldst fear to do in mine.
\end{quotation}

Bernard keeps that affection anchored squarely in God: like stars shining beneath the sun, ``the honour due to God must prevail over and in a sense exclude the honour claimed by others, so that He alone is not only honoured more than any other, but is honoured in all others.'' Urging his monks to ``make the angels of God your familiar friends; frequent their society by constant remembrance and fervent prayer,'' he adds in Sermon~XIII that angelic hands bear the soul safely along ``that last long journey'' past every demonic snare.

In Carmelite mysticism, St.~Teresa of \'Avila records in her \emph{Life} (ch.~XXIX, \S\S16--17) the extraordinary bodily vision of the transverberation---noting that she usually perceived angels only by intellectual vision, and remarking with theological humility that angels ``never tell me'' their names:

\begin{quotation}
He was not large, but small of stature, and most beautiful---his face burning, as if he were one of the highest angels, who seem to be all of fire: they must be those whom we call cherubim \dots\ I saw in his hand a long spear of gold, and at the iron's point there seemed to be a little fire. He appeared to me to be thrusting it at times into my heart, and to pierce my very entrails; when he drew it out, he seemed to draw them out also, and to leave me all on fire with a great love of God. The pain was so great, that it made me moan; and yet so surpassing was the sweetness of this excessive pain, that I could not wish to be rid of it. The soul is satisfied now with nothing less than God. The pain is not bodily, but spiritual; though the body has its share in it, even a large one.
\end{quotation}

(Teresa's own Escorial autograph reads ``cherubim,'' whereas earlier printed editions adopted her confessor Domingo B\'a\~nez's marginal suggestion ``seraphim.'')

St.~Francis de Sales translates angelic communion into everyday lay devotion in the \emph{Introduction to the Devout Life} (Part~II, ch.~16): ``Seek to be familiar with the Angels; learn to realise that they are continually present, although invisible. Specially love and revere the Guardian Angel of the Diocese in which you live, those of the friends who surround you, and your own. Commune with them frequently, join in their songs of praise, and seek their protection and help in all you do, spiritual or temporal.'' Francis illustrates this counsel with St.~Peter Faber, first companion of St.~Ignatius, who when traveling through troubled territories in Germany made a practice of greeting the guardian angels of every parish he entered, ``whose help had often preserved him from danger, and softened hearts to receive the faith.''

St.~Alphonsus Liguori draws an urgent lesson for perseverance from the fall of the rebel angels in \emph{Preparation for Death}. Quoting St.~Gelasius---``receiving the grace of God in vain, they could not persevere, because they did not pray''---Alphonsus asks: if pure spirits free from flesh, world, and tempter still fell for failing to pray, how could weak human beings persevere till death without continual prayer? In his method of mental prayer (\emph{The Incarnation, Birth, and Infancy of Jesus Christ}), he builds angelic devotion into daily discipline: ``say a Hail Mary to the divine Mother, and a Glory be to the Father, etc., in honor of St.~Joseph, of your guardian angel, and of your holy patron.''

St.~Th\'er\`ese of Lisieux unites childlike intimacy with austere doctrinal sobriety. As a child at bedtime, she relates in \emph{Story of a Soul}, she asked Pauline every evening, ``Have I been good to-day? Is God pleased with me? Will the Angels watch over me?'' On pilgrimage in Naples, when carriage horses bolted down a steep hill, ``it is solely to our Guardian Angels that I attribute our safe return.'' Preparing for Holy Communion, she begged Our Lady to make her soul a tabernacle and then invited ``all the Angels and Saints to come and sing canticles of love.'' Near death (\emph{Counsels and Reminiscences}), she grounded her famous promise of heavenly intercession directly in angelic ministry: ``I WILL SPEND MY HEAVEN IN DOING GOOD UPON EARTH. Nor is this impossible, since from the very heart of the Beatific Vision, the Angels keep watch over us.'' Yet when a sister tried to comfort her with pictures of white-robed angels carrying her soul to heaven, Th\'er\`ese spoke with Thomistic clarity: ``Fancies like those,'' she answered, ``do not help me, and my soul can only feed upon truth. God and His Angels are pure spirits. No human eye can see them as they really are. That is why I have never asked extraordinary favours. I prefer to await the Eternal Vision.''

\subsection{One City Under Christ}

Created out of nothing at the dawn of time, pure immaterial intellects without bodies, the holy angels stand in nine radiant choirs before the face of the Triune God---yet they are neither remote gods nor rivals to mankind. Every truth of Catholic angelology converges upon a single commonwealth. In nature, the least angel surpasses every mortal body; in grace, redeemed human beings are called to stand beside the Seraphim, while Christ in his assumed humanity reigns as Head of angels and men alike, with Mary his Mother exalted as Queen above every choir.

Between heaven and earth stretches no unbridgeable chasm, for the angels ascend and descend upon the Son of Man (John 1:51), illuminating minds, guarding souls from birth to death, restraining the malice of the fallen spirits, and carrying the Church's sacrifice to the altar on high. Even now, the pilgrim believer lives inside the society described to the Hebrews: ``But you are come to mount Sion and to the city of the living God, the heavenly Jerusalem, and to the company of many thousands of angels, and to the church of the firstborn who are written in the heavens, and to God the judge of all, and to the spirits of the just made perfect, and to Jesus the mediator of the new testament'' (Heb 12:22--24).
